\documentclass[letterpaper, 10 pt, conference]{ieeeconf}

\IEEEoverridecommandlockouts
\overrideIEEEmargins

\usepackage[final]{graphicx}
\usepackage{hyperref}
\usepackage{amsmath}
\usepackage{amsfonts}
\usepackage{float}
\graphicspath{{images/}}
\usepackage{mathtools}
\usepackage{cite}
\usepackage{tabularx}
\usepackage{multirow}
\usepackage{booktabs}
\usepackage{adjustbox}
\usepackage{bm}
\usepackage{hhline,array}
\usepackage[dvipsnames]{xcolor}

% Shared macros across all three papers. Kept minimal: these papers are
% not math-heavy; most quantities are named directly in text/equations.
\newcommand{\surfz}{z_{\text{surf}}}
\newcommand{\topz}{z_{\text{top}}}
\newcommand{\height}{h}
\newcommand{\descent}{d}
\newcommand{\graspw}{w_g}
\newcommand{\tcpoff}{\Delta z_{\text{tcp}}}
\newcommand{\code}[1]{\texttt{#1}}
\newcommand{\fb}[1]{\textcolor{orange}{#1}}
% Marks text changed in this revision round (reviewer-requested edits) so
% readers can see exactly what moved. Strip \rev{...} -> {...} once the
% revision is accepted.
\newcommand{\rev}[1]{\textcolor{red}{#1}}


\begin{document}

\title{\LARGE \bf Don't Trust the Model's Arithmetic: Deterministic Geometric Grounding for LLM/VLM-Driven Manipulation}

\author{Matteo Nini$^{1}$%
	\thanks{$^{1}$The authors are with [Department, University]. E-mail: {\tt\small matteo.nini@[domain]}.}%
	% TODO(authors): add co-authors and affiliation block, matching the
	% format used in prior lab submissions (see old_paper_extracted/root.tex).
}

\maketitle

\begin{abstract}
Large language model (LLM) and vision-language model (VLM) based robotic
planners are increasingly validated in simulation, where an object's
position, size, and grasp geometry are either exactly known or scripted.
We report on a real UR5cb manipulator deployment where these planners are
routinely asked to compute grasp and release geometry themselves --
release height, target stacking height, grasp width -- and show that this
model-computed arithmetic, not high-level task reasoning, is the dominant
source of planning failure on physical hardware. Across six development
sessions we diagnose a recurring family of geometry-grounding bugs (e.g.\
a systematic release-height overshoot from a top-vs-mid-body grasp
mismatch, and a double-counted stacking height in a hybrid vision-language
pipeline), each invisible to a symbolic or scripted simulator because
none of these quantities are ever actually measured there. We fix five of
these deterministically, removing geometric arithmetic from the model
entirely and computing it from paired pick/release actions and measured
depth; the sixth, a multi-object placement collision that recurs when
sequential releases share a target zone, we diagnose but do not yet close,
and report as the dominant residual failure. On a 120-trial real-hardware
benchmark comparing an LLM planner against a camera-grounded VLM planner
using a single unified model for both, we show this residual failure
distribution occurs in \textbf{both} arms regardless of grounding
modality, and argue that symbolic-simulation benchmarks of LLM-based
robot planners -- including a directly comparable prior taxonomy-driven
study -- structurally cannot observe this failure class, and that a
geometry-grounding audit belongs in the evaluation protocol for any
LLM/VLM planner headed to real hardware.
\end{abstract}

\section{Introduction}
\label{sec:introduction}
Large language models (LLMs) and vision-language models (VLMs) are
increasingly used as high-level planners for robotic manipulation,
translating a natural-language instruction and a scene observation into a
sequence of executable skills~\cite{kim2024survey, wang2025robotics,
ahn2022saycan, huang2022zeroshot}. A large and fast-growing body of work
validates these planners in simulation: scripted physics engines such as
CoppeliaSim~\cite{rohmer2013coppeliasim} or MuJoCo~\cite{todorov2012mujoco},
or, further still, purely symbolic environments in which the ``scene'' is a
hand-authored text description and perception is a proxy that returns
whatever the experiment designer wrote down~\cite{favaliRAS}. Symbolic and
scripted validation is valuable -- it is fast, reproducible, and lets a
researcher isolate the reasoning strategy from perception and control. But
it has a structural blind spot: an object's height, its grasp width, and
the vertical distance a gripper must travel to release it safely are never
actually \emph{measured} in these environments. They are simply asserted,
correctly, by construction.

On real hardware, none of these quantities are given. They must be
computed, from noisy depth, from a calibrated but imperfect camera, or from
a hand-authored scene file that may or may not match the physical table it
describes. Over six development sessions deploying an LLM/VLM planning
stack on a real UR5cb manipulator with an OnRobot RG2 gripper and a
calibrated RGB-D camera, we observed a recurring and costly pattern: when a
planning prompt asks the model to compute a piece of manipulation geometry
itself -- a release height, a stacking height, a grasp width -- that
computation is often wrong, in ways that are silent until the arm
physically drops, crushes, or misses an object. The failures are not
random. They cluster around a single, identifiable root cause: the model
is asked to reason about a physical quantity that a sensor, not language,
should determine.

This paper makes that observation precise and actionable. We diagnose six
recurring geometry-grounding bugs from production hardware logs (\S
\ref{sec:taxonomy}), show that they share a common structural cause across
two independently-implemented planning pipelines (a static-scene LLM
planner and a camera-grounded VLM planner), and close five of them by
removing geometric arithmetic from the model's output entirely, replacing
it with a deterministic grounding layer that computes the same quantities
from paired pick/release actions and measured depth
(\S\ref{sec:grounding-layer}). The sixth -- a multi-object placement
collision -- we diagnose but have not yet closed, and it turns out to be
the dominant one: a 120-trial real-hardware benchmark (\S\ref{sec:evaluation})
shows it as the residual failure distribution once the other five are
fixed, occurring in \emph{both} the LLM and VLM pipelines regardless of
which one grounded the scene -- direct evidence that the failure tracks
the geometry-grounding architecture, not the choice of grounding
modality.

Our central argument is that this failure class is not a peculiarity of
our codebase, but a structural consequence of validating LLM/VLM robot
planners in simulation. We make this argument concretely against a direct,
comparable prior work: Favali et al.~\cite{favaliRAS} propose a rigorous
taxonomy of task and environment complexity and a benchmarking framework
for LLM-based robotic planning, evaluated in a symbolic simulator where
ground truth is always exactly known. Their framework is, by design,
incapable of producing the failure mode this paper documents, because
their simulated agent never computes a grasp height that a physical sensor
could contradict. We do not consider this a weakness of their study --
their methodology is sound for the question it asks -- but we argue that a
geometry-grounding audit of the kind we perform here is a necessary
complement to any symbolic or scripted benchmark before an LLM/VLM planner
is deployed on real hardware.

The contributions of this paper are threefold:
\begin{itemize}[nosep]
  \item A taxonomy of six recurring real-hardware geometry-grounding
    bugs, diagnosed from production debug logs across six development
    sessions on a real manipulator (\S\ref{sec:taxonomy}).
  \item A deterministic grounding layer that closes five of the six,
    removing model-computed geometric arithmetic from the
    planning-critical path, together with an evaluation of its effect
    (\S\ref{sec:grounding-layer}).
  \item A 120-trial real-hardware benchmark, methodologically anchored to
    a directly comparable symbolic-simulation study~\cite{favaliRAS} and
    using a single unified model for both grounding modalities, showing
    that the sixth, still-open bug -- multi-object placement collision --
    is the dominant residual failure in \emph{both} arms, and what that
    implies for evaluating LLM/VLM robot planners before hardware
    deployment (\S\ref{sec:evaluation}).
\end{itemize}


\section{Related Work}
\label{sec:related-work}
\textbf{LLM/VLM planning for robotics.} Foundation models have been used to
ground natural-language instructions in robotic
affordances~\cite{ahn2022saycan}, to act as zero-shot
planners~\cite{huang2022zeroshot}, and to generate feasible motion plans
directly from instructions~\cite{lin2023text2motion}. A range of reasoning
strategies has been proposed on top of these capabilities: reactive,
single-step prompting; interleaved reasoning-and-acting
(ReAct)~\cite{yao2022react}; iterative self-critique
(Self-Refine)~\cite{madaan2023selfrefine}, including a long-horizon
sequential-planning variant validated at ICRA~\cite{zhou2024isrllm};
search-based deliberation (Tree-of-Thoughts)~\cite{yao2023tot}; and
ensemble voting over independently sampled plans
(chain-of-thought with self-consistency)~\cite{wei2022cot,
wang2022selfconsistency}. Recent surveys catalogue this space in
detail~\cite{kim2024survey, wang2025robotics, huang2024planningsurvey,
berti2025emergent, billard2025roadmap}. Our contribution is orthogonal to
this line of work: we do not propose a new reasoning strategy, and hold
the reasoning strategy fixed in our evaluation (\S\ref{sec:evaluation}).
Instead, we show that the geometric grounding surrounding \emph{any} of
these strategies -- how a chosen action's numeric parameters are computed
-- is a separate failure axis that persists regardless of which reasoning
strategy produced the action.

\textbf{End-to-end and vision-language-action approaches.} An alternative
architecture collapses perception, reasoning, and control into a single
model, as in vision-language-action (VLA) systems~\cite{brohan2023rt2,
driess2023palme} or code-generating policies that couple perception
directly to control primitives~\cite{liang2023codeaspolicies,
huang2023voxposer, huang2022innermonologue, singh2023progprompt}. These
approaches can, in principle, learn to compensate for geometric grounding
errors end-to-end, at the cost of the diagnosability that a decoupled
Sense-Plan-Act pipeline provides -- when an end-to-end model fails, it is
difficult to determine whether the failure was a misperception, a
reasoning error, or a motor-control error~\cite{favaliRAS}. We deliberately
retain a decoupled architecture, precisely because it lets us isolate and
name the geometry-grounding failure class this paper documents; a fully
end-to-end system would not permit the diagnosis in \S\ref{sec:taxonomy}
at all, only its symptoms.

\textbf{Simulation and sim-to-real transfer.} The sim-to-real gap is a
long-studied concern in robot learning, primarily for learned low-level
control policies~\cite{zhao2020simtoreal}. Physics simulators such as
Gazebo~\cite{koenig2004gazebo}, MuJoCo~\cite{todorov2012mujoco}, and
CoppeliaSim~\cite{rohmer2013coppeliasim} narrow this gap for control by
modeling contact dynamics and actuation limits~\cite{collins2021physics}.
The gap we document is different in kind: it does not concern whether a
\emph{learned policy} generalizes from simulated to real dynamics, but
whether a \emph{language-model-computed} geometric quantity -- one that a
symbolic benchmark never forces the model to compute against a
contradicting sensor -- is trustworthy at all. To our knowledge this
failure class has not been previously isolated and named for LLM/VLM robot
planners specifically.

\textbf{Direct comparison.} The closest prior work to this paper is
Favali et al.~\cite{favaliRAS}, which introduces a task/environment
complexity taxonomy (task length, goal specificity, affordance
perception; affordance density and four classes of dynamic disturbance)
and TS/TSR/AETS metrics for benchmarking LLM-based robotic planning
across six reasoning strategies, evaluated in a symbolic simulator with
exact ground truth. We adopt the same evaluation philosophy and, for
\S\ref{sec:evaluation}, the same metric family, but on real hardware,
where ground truth for safety and success is not automatically available
(\S\ref{sec:evaluation} details the resulting semi-automatic
adaptation). A companion systems paper from the same group,
RoboReason~Lab~\cite{favaliRoboReasonLab}, provides a modular simulation
testbed for the same class of reasoning strategies in CoppeliaSim; like
\cite{favaliRAS}, it never grounds a plan in measured depth or a
calibrated camera. Both are simulation-only by design, and both are
structurally unable to produce the geometry-grounding failures this paper
diagnoses, because neither ever asks a model to compute a quantity a
physical sensor could contradict.


\section{Geometry-Grounding Failure Taxonomy}
\label{sec:taxonomy}
\input{sections/03_geometry_taxonomy}

\section{Deterministic Grounding Layer}
\label{sec:grounding-layer}
\input{sections/04_grounding_layer}

\section{Real-Hardware Evaluation}
\label{sec:evaluation}
\input{sections/05_evaluation}

\section{Discussion and Limitations}
\label{sec:discussion}
\textbf{Implication for evaluation practice.} If the argument of this
paper is correct, symbolic and scripted simulation is necessary but not
sufficient for validating an LLM/VLM robot planner: it can validate
reasoning strategy, task decomposition, and high-level correctness, but by
construction cannot exercise the geometry-grounding failure class of
\S\ref{sec:taxonomy}, because it never asks the model to compute a
quantity a physical sensor could contradict. We suggest that a
geometry-grounding audit -- identifying every numeric quantity a planning
prompt asks the model to compute itself, and asking whether a sensor or a
paired action already determines that quantity more reliably -- belongs
alongside task-complexity benchmarking as a standard step before hardware
deployment, independent of which reasoning strategy or foundation model is
used.

\textbf{What the missing ablation would add.} \S\ref{sec:validation} is
explicit that our deterministic grounding layer is validated
retrospectively, against replayed real-hardware failures and a unit
regression suite, rather than through a live controlled A/B comparison
on hardware. A live ablation -- re-running the task conditions that
originally surfaced these bugs with the grounding layer enabled versus
disabled -- would let us report a quantitative safety/success delta
attributable to the fix specifically, rather than inferring its
necessity from the failures it was designed to close. Two scopes are
designed but neither is run at the time of writing: a restricted version
(LLM-only, the two release-geometry fixes specifically, roughly half a
day including trials) and a complete one (both modalities, the full
taxonomy including the safety-critical TCP-offset/grasp-width path,
roughly two days). We flag the restricted version as the most direct
next step, since it targets the mechanism the evaluation of
\S\ref{sec:evaluation} actually implicates.

\textbf{Case 6 remains the largest open item.} Everything else in this
paper is either closed (five of six bugs) or bounded by a controlled
comparison we can eventually run (the ablation above). Case 6 is
neither: it is a trajectory-level, multi-body contact problem, and
\S\ref{sec:zone-spacing}'s static-spacing mitigation, while a genuine
improvement over an unbounded layout, does not touch the dynamic half of
the problem at all. Closing it plausibly needs one of (a) a
motion-planning-level check of the approach trajectory against
already-placed objects' known footprints, not just their final target
positions, or (b) a settle-time verification step after each release,
before the next pick begins, so a displaced neighbor is caught and
corrected rather than compounding into the next release. We have not
implemented or evaluated either.

\textbf{Generalizability.} Our evidence comes from one manipulator (UR5cb),
one gripper (OnRobot RG2), one camera rig, and one reasoning strategy.
Deliberately, it comes from a \emph{single} underlying model
(\code{kimi-k2.6}) used for both the LLM and VLM arm, rather than one
model per modality -- an earlier configuration with modality-specific
models is what originally produced the now-non-replicating
value-mapping-reasoning finding of \S\ref{sec:computation-not-perception},
which in hindsight looks like it was at least partly a model-identity
confound rather than a pure grounding-modality effect. We do not claim
that the specific six bugs of Table~\ref{tab:taxonomy} recur verbatim on
different hardware; we claim that the \emph{structural pattern} --
geometry computed by a model rather than measured or derived from a
paired action, and multiple real bodies sharing space over time -- is a
property of the architecture (asking an LLM/VLM to output numeric
geometry at all, and placing more than one object in a shared zone), not
of this particular robot, and should be audited for on any comparable
deployment.

\textbf{Grasp width remains unverified.} Of the quantities in
\S\ref{sec:grounding-layer}, grasp width is the one our system still
cannot cross-check against depth in every case: a single top-down depth
sample determines an object's height unambiguously, but not its
left-right extent. Our best-effort edge-inset sampling
(\S\ref{sec:grounding-layer}) reduces, but does not eliminate, this gap;
a depth-informed rectification of the camera view before grounding is a
natural next step we have not yet completed.

\textbf{Reliability of operator-supplied labels.} Our semi-automatic
evaluation protocol (\S\ref{sec:evaluation}) depends on a single
operator's real-time judgment of safety and sub-task completion for each
trial. We did not measure inter-rater reliability with a second, blinded
annotator on a subsample of trials, which would strengthen confidence in
the reported TS/TSR numbers independent of the geometry-grounding
argument this paper makes. This is not a hypothetical risk: an early pass
over this cycle's data marked every trial safe regardless of outcome,
caught only by cross-checking the safety label against the free-text
operator note recorded for the same trial, and corrected before the
numbers in \S\ref{sec:evaluation} were computed. We report the incident
as a concrete illustration of exactly the reliability concern this
paragraph raises, not as a suggestion that it recurs undetected
elsewhere in this dataset.


\section{Conclusion}
\label{sec:conclusion}
We diagnosed a recurring class of geometry-grounding failures in LLM/VLM
robot planning, drawn from six sessions of real-hardware deployment, and
showed that its structural cause -- a model computing a physical quantity
that a sensor, a paired action, or real multi-body contact already
determines -- is invisible to symbolic and scripted simulation by
construction, including a directly comparable prior taxonomy-driven
benchmark. Replacing model-computed geometry with a deterministic
grounding layer closes five of the six bugs we diagnosed; a 120-trial
real-hardware evaluation with that layer in place, using a single unified
model for both grounding modalities, shows the sixth -- a multi-object
placement collision -- as the dominant residual failure, occurring in
both the LLM and VLM pipelines regardless of which one grounded the
scene. We argue that a geometry-grounding audit is a necessary complement
to task-complexity benchmarking for any LLM/VLM planner headed to
physical hardware, and leave closing the sixth bug's dynamic,
trajectory-level half, a controlled hardware ablation isolating the
grounding layer's quantitative effect, and a depth-verified treatment of
grasp width, as the most direct next steps.


% NOTE: IEEEtran.bst is not installed in this local TeX Live; using ieeetr
% as a local stand-in. Overleaf and the official IEEE/ICRA LaTeX template
% both bundle IEEEtran.bst -- switch back to \bibliographystyle{IEEEtran}
% for the final submission build.
\bibliographystyle{ieeetr}
\bibliography{bibliography}

\end{document}
